\lecture{18}{Масштабированный градиентный поток}

\sourcevideo
    {https://www.youtube.com/playlist?list=PLOzRYVm0a65fhKDS8cYIC7YCuVGJ4FOJx}
    {Week 5: Lecture 18: Rescaled Gradient Flow}

Обычный градиентный поток замедляется около стационарной точки,
поскольку норма градиента стремится к нулю. Положительное скалярное
масштабирование сохраняет геометрические траектории непрерывной
системы, но изменяет скорость их прохождения. Подходящий степенной
множитель позволяет получить точное достижение оптимума сначала за
конечное, а затем за фиксированное время.

\section{Масштабирование градиентного потока}

Рассмотрим задачу
\[
    \min_{x\in\RR^n}f(x)
\]
и обычный градиентный поток
\[
    \dot x=-\nabla f(x).
\]
Его равновесия совпадают со стационарными точками функции. Если
\(f\) сильно выпукла, стационарная точка единственна и равна
глобальному минимуму \(x^\star\). Скорость обычного потока равна
\[
    \lVert\dot x\rVert=\lVert\nabla f(x)\rVert,
\]
поэтому движение замедляется при приближении к решению.

Пусть \(p>2\) и
\[
    \rho_p=\frac{p-2}{p-1}\in(0,1).
\]
Масштабированный поток определяется как
\[
    \dot x
    =
    -\frac{\nabla f(x)}
    {\lVert\nabla f(x)\rVert^{\rho_p}}
\]
при \(\nabla f(x)\ne0\), а в стационарных точках полагается
\(\dot x=0\). Его скорость равна
\[
    \lVert\dot x\rVert
    =
    \lVert\nabla f(x)\rVert^{1-\rho_p}
    =
    \lVert\nabla f(x)\rVert^{1/(p-1)}.
\]
Поскольку \(1/(p-1)<1\), скорость около минимума убывает медленнее,
чем у обычного градиентного потока.

Вне стационарных точек масштабированное поле имеет вид
\[
    \widetilde F(x)=a(x)F(x),
    \qquad
    a(x)=\lVert\nabla f(x)\rVert^{-\rho_p}>0,
\]
где \(F(x)=-\nabla f(x)\). Поэтому оба поля задают одинаковые
непараметризованные фазовые траектории и различаются только
перепараметризацией времени. После дискретизации это свойство в общем
случае теряется.

Величина правой части около равновесия имеет порядок
\[
    \lVert\nabla f(x)\rVert^{1/(p-1)}.
\]
Показатель меньше единицы, поэтому поле обычно гёльдерово, но не
локально липшицево в стационарной точке. Именно такая потеря
липшицевости допускает точное достижение равновесия за конечное время.

\section{Конечно-временной критерий Ляпунова}

\begin{definition}[Конечно-временная устойчивость]
Равновесие \(x^\star\) называется конечно-временно устойчивым, если
оно устойчиво по Ляпунову и для каждого допустимого начального
состояния \(x_0\) существует число \(T(x_0)<\infty\), такое что
\[
    x(t)=x^\star
\]
для всех \(t\ge T(x_0)\). Число \(T(x_0)\) называется временем
установления.
\end{definition}

\begin{theorem}[Конечно-временной критерий]
Пусть \(V\) положительно определена относительно \(x^\star\) и вдоль
траекторий выполняется
\[
    \dot V\le-cV^\alpha,
    \qquad
    c>0,
    \qquad
    0<\alpha<1.
\]
Тогда равновесие достигается за конечное время, причём
\[
    T(x_0)
    \le
    \frac{V(x_0)^{1-\alpha}}
    {c(1-\alpha)}.
\]
\end{theorem}

\begin{proof}
Пока \(V(t)>0\), имеем
\[
    V(t)^{-\alpha}\dot V(t)\le-c.
\]
Следовательно,
\[
    \frac{\mathrm d}{\mathrm dt}V(t)^{1-\alpha}
    \le
    -c(1-\alpha).
\]
После интегрирования получаем
\[
    V(t)^{1-\alpha}
    \le
    V(x_0)^{1-\alpha}-c(1-\alpha)t.
\]
Правая часть обращается в нуль не позднее указанного момента. Так как
\(V\ge0\), после этого должно выполняться \(V(t)=0\), а положительная
определённость даёт \(x(t)=x^\star\).
\end{proof}

Отличие от экспоненциального условия
\[
    \dot V\le-cV
\]
состоит в показателе \(\alpha<1\). При \(0<V<1\) выполняется
\(V^\alpha>V\), поэтому убывание около равновесия достаточно быстро
для точного достижения нуля.

\section{Конечное время масштабированного потока}

Предположим, что \(f\in C^2(\RR^n)\) и
\[
    \nabla^2f(x)\succeq\mu I,
    \qquad
    \mu>0.
\]
Тогда функция сильно выпукла, минимизатор \(x^\star\) единственен и
\[
    \nabla f(x)=0
    \quad\Longleftrightarrow\quad
    x=x^\star.
\]
Выберем функцию Ляпунова
\[
    V(x)=\frac12\lVert\nabla f(x)\rVert^2.
\]
Она положительно определена относительно \(x^\star\). Вдоль
масштабированного потока
\[
\begin{aligned}
    \dot V
    &=
    -\frac{
        \nabla f(x)^{\mathsf T}
        \nabla^2f(x)
        \nabla f(x)
    }{
        \lVert\nabla f(x)\rVert^{\rho_p}
    }
    \\
    &\le
    -\mu
    \lVert\nabla f(x)\rVert^{2-\rho_p}.
\end{aligned}
\]
Поскольку
\[
    2-\rho_p=\frac{p}{p-1},
\]
и \(\lVert\nabla f(x)\rVert^2=2V(x)\), получаем
\[
    \dot V
    \le
    -\mu2^{\alpha_p}V^{\alpha_p},
    \qquad
    \alpha_p=\frac{p}{2(p-1)}.
\]
При \(p>2\)
\[
    \frac12<\alpha_p<1,
\]
поэтому применим конечно-временной критерий.

\begin{theorem}[Конечно-временная сходимость]
Пусть \(f\in C^2(\RR^n)\) является \(\mu\)-сильно выпуклой и
\(p>2\). Тогда равновесие \(x^\star\) масштабированного потока
конечно-временно устойчиво, причём
\[
    T(x_0)
    \le
    \frac{p-1}{\mu(p-2)}
    \lVert\nabla f(x_0)\rVert^{\frac{p-2}{p-1}}.
\]
\end{theorem}

Время конечно, но зависит от начального состояния. В конечно-временном
случае после момента \(T\) выполняется \(V(t)=0\). При сильной
выпуклости это непосредственно означает
\[
    \nabla f(x(t))=0
    \quad\Longleftrightarrow\quad
    x(t)=x^\star,
\]
поэтому дополнительная двусторонняя оценка между \(V\) и расстоянием
до решения не требуется.

В лекции также упоминаются полиномиальные оценки вида
\[
    f(x(t))-f^\star=O(t^{-p}).
\]
Их нельзя напрямую сравнивать с полученным конечно-временным
результатом: такие оценки могут относиться к более широкому классу
выпуклых функций, тогда как приведённый вывод использует равномерную
нижнюю границу Гессиана.

\section{Фиксированное время и два режима}

\begin{definition}[Фиксированно-временная устойчивость]
Равновесие называется фиксированно-временно устойчивым, если оно
конечно-временно устойчиво и существует число \(T_{\max}<\infty\), не
зависящее от начального состояния, для которого
\[
    T(x_0)\le T_{\max}
\]
для всех допустимых \(x_0\).
\end{definition}

\begin{theorem}[Фиксированно-временной критерий]
Пусть \(V\) положительно определена относительно \(x^\star\) и
\[
    \dot V
    \le
    -c_1V^{\alpha_1}
    -c_2V^{\alpha_2},
\]
где
\[
    c_1,c_2>0,
    \qquad
    0<\alpha_1<1<\alpha_2.
\]
Тогда
\[
    T(x_0)
    \le
    T_{\max}
    =
    \frac1{c_1(1-\alpha_1)}
    +
    \frac1{c_2(\alpha_2-1)}.
\]
\end{theorem}

Доказательство разбивает движение на два этапа. При \(V\ge1\)
оставляем только второе слагаемое:
\[
    \dot V\le-c_2V^{\alpha_2}.
\]
Интегрирование от \(V(x_0)\) до единицы даёт время не больше
\[
    \frac1{c_2(\alpha_2-1)}.
\]
После входа в область \(V\le1\) используется
\[
    \dot V\le-c_1V^{\alpha_1},
\]
и дополнительное время до нуля не превосходит
\[
    \frac1{c_1(1-\alpha_1)}.
\]
Сумма не содержит начального значения. Степень выше единицы
контролирует движение вдали от решения, а степень ниже единицы
обеспечивает точное достижение около него.

\section{Двухрежимный масштабированный поток}

Выберем
\[
    p>2,
    \qquad
    1<q<2,
\]
и положим
\[
    \rho_p=\frac{p-2}{p-1},
    \qquad
    \rho_q=\frac{q-2}{q-1}.
\]
Тогда \(0<\rho_p<1\), а \(\rho_q<0\). Рассмотрим систему
\[
\begin{aligned}
    \dot x
    ={}&
    -\frac{\nabla f(x)}
    {\lVert\nabla f(x)\rVert^{\rho_p}}
    \\
    &-
    \frac{\nabla f(x)}
    {\lVert\nabla f(x)\rVert^{\rho_q}}
\end{aligned}
\]
при \(\nabla f(x)\ne0\), а в стационарной точке положим
\(\dot x=0\).

Первое слагаемое имеет величину
\[
    \lVert\nabla f(x)\rVert^{1/(p-1)}
\]
с показателем меньше единицы и действует около решения. Второе имеет
величину
\[
    \lVert\nabla f(x)\rVert^{1/(q-1)}
\]
с показателем больше единицы и доминирует вдали от решения.

Для
\[
    V(x)=\frac12\lVert\nabla f(x)\rVert^2
\]
и \(\nabla^2f(x)\succeq\mu I\) получаем
\[
    \dot V
    \le
    -\mu2^{\alpha_1}V^{\alpha_1}
    -\mu2^{\alpha_2}V^{\alpha_2},
\]
где
\[
    \alpha_1=\frac{p}{2(p-1)}<1,
    \qquad
    \alpha_2=\frac{q}{2(q-1)}>1.
\]

\begin{theorem}[Фиксированное время]
Пусть \(f\in C^2(\RR^n)\) является \(\mu\)-сильно выпуклой,
\(p>2\) и \(1<q<2\). Тогда оптимальное равновесие двухрежимного
потока фиксированно-временно устойчиво и
\[
\begin{aligned}
    T(x_0)
    \le T_{\max}
    ={}&
    \frac1{
        \mu2^{\alpha_1}(1-\alpha_1)
    }
    \\
    &+
    \frac1{
        \mu2^{\alpha_2}(\alpha_2-1)
    }.
\end{aligned}
\]
Правая часть не зависит от \(x_0\).
\end{theorem}

Эта конструкция иллюстрирует проектирование динамики по требуемому
неравенству Ляпунова. Сначала задаются желаемые степени
\(V^{\alpha_1}\) и \(V^{\alpha_2}\), затем выбирается функция
Ляпунова, после чего правая часть системы строится так, чтобы её
производная содержала эти степени.

\section{Регуляризация, дискретизация и стоимость}

В численной реализации знаменатель часто заменяют на
\[
    \lVert\nabla f(x)\rVert+\varepsilon,
    \qquad
    \varepsilon>0.
\]
Например,
\[
    \dot x
    =
    -\frac{\nabla f(x)}
    {\bigl(\lVert\nabla f(x)\rVert+\varepsilon\bigr)^{\rho_p}}.
\]
Это устраняет деление на нуль, но около оптимума поток ведёт себя как
обычный градиентный поток с постоянным масштабом
\(\varepsilon^{-\rho_p}\). Поэтому при фиксированном
\(\varepsilon>0\) точная конечно-временная гарантия обычно теряется;
остаётся асимптотическая сходимость или быстрое попадание в малую
окрестность.

Явная дискретизация одночленного потока имеет вид
\[
    x^{k+1}
    =
    x^k
    -\eta
    \frac{\nabla f(x^k)}
    {\lVert\nabla f(x^k)\rVert^{\rho_p}}.
\]
Для двухрежимной системы к ней добавляется аналогичный член с
\(\rho_q\). Непрерывная конечно- или фиксированно-временная
устойчивость не переносится автоматически на эту рекурсию. Около
решения возможны перескакивание через минимум, колебания и высокая
чувствительность к округлению. Поэтому шаг, регуляризация и критерий
остановки требуют отдельного анализа.

По сравнению с обычным градиентным методом дополнительно вычисляются
норма градиента и одна или две скалярные степени. Для вектора
размерности \(n\) вычисление нормы и масштабирование требуют
\(O(n)\) операций, тогда как основная стоимость обычно связана с
самим градиентом. В распределённой системе глобальная норма может
потребовать отдельного агрегирования.

Теоремы предполагают точный непрерывный поток, точный градиент,
\(f\in C^2\) и равномерную сильную выпуклость. Стохастический шум,
задержки, коммуникационные ошибки и конечный шаг могут уничтожить
точное достижение решения. Эмпирические сравнения с Adam и RMSProp зависят от задачи и параметров и сами по себе не устанавливают универсального превосходства одного метода.
